\documentclass{article}
\usepackage[UTF8]{ctex}
\usepackage{amsmath,amssymb,amsfonts}
\usepackage{enumitem}
\usepackage{geometry}
\geometry{a4paper, margin=1in}
% === 错题本排版优化（不参与题目标识） ===
\setlength{\parindent}{0pt}
\setlength{\parskip}{0.35em plus 0.10em minus 0.08em}
\setlength{\emergencystretch}{3em}
\setlist{itemsep=0.25em, topsep=0.40em, parsep=0pt}
\allowdisplaybreaks[2]
% === 错题本排版优化结束 ===
\begin{document}
\section*{南京航空航天大学2025年数学分析真题}
\section*{99．南京航空航天大学2025年数学分析试卷}

\subsection{题目1}
一．求下列极限和积分。

\begin{itemize}
  \item[（1）] \(\lim _{x \rightarrow 0}\left(\frac{e^{x}+e^{2 x}+e^{3 x}}{3}\right)^{\frac{1}{x}}\) ．
  \item[（2）] \(\int_{0}^{1} \frac{e^{x}(1-x)^{2}}{\left(1+x^{2}\right)^{2}} \mathrm{~d} x\) ．
\end{itemize}

\subsection{题目2}
二．设 \(f(x)\) 在 \([a,+\infty)\) 上有定义，证明： \(\lim _{x \rightarrow+\infty} f(x)=A\) 的充要条件是对任意包含在 \([a,+\infty)\) 上的单调递增数列 \(\left\{x_{n}\right\}\) ，如果 \(\lim _{n \rightarrow \infty} x_{n}=+\infty\) ，就有 \(\lim _{n \rightarrow \infty} f\left(x_{n}\right)=A\) 。

\subsection{题目3}
三．解答如下问题：\\
续．（1）设 \(f(x)\) 在 \((a, b)\) 上连续，且 \(\lim _{x \rightarrow a^{+}} f(x)\) 与 \(\lim _{x \rightarrow b^{-}} f(x)\) 均存在且有限，证明：\(f(x)\) 在 \([a, b]\) 上一致连\\
（2）证明：\(f(x)=\frac{|\sin x|}{x}\) 在（0，1）与（－1，0）上一致连续，但在 \(0<|x|<1\) 上不一致连续．

\subsection{题目4}
四．设 \(0<a<b, f(x)\) 在 \([a, b]\) 上连续，在 \((a, b)\) 上可导．

\begin{itemize}
  \item[（1）] 证明：存在 \(\xi \in(a, b)\) ，使得 \(f(b)-f(a)=\xi f^{\prime}(\xi) \ln \frac{b}{a}\) ．
  \item[（2）] 证明： \(\lim _{n \rightarrow \infty} n(\sqrt[n]{x}-1)=\ln x\) ．
\end{itemize}

\subsection{题目5}
五．设 \(f(x)\) 在 \([a, b]\) 上二阶可导，且 \(f^{\prime \prime}(x) \geq 0, f(x) \leq 0\) ．证明：

\[
f(x) \geq \frac{2}{b-a} \int_{a}^{b} f(x) \mathrm{d} x, x \in[a, b]
\]

\subsection{题目6}
六．设

\[
f_{n}(x)=\sum_{k=1}^{n} \frac{1}{n} \cos \left(x+\frac{k}{n}\right), n=1,2, \cdots .
\]

求 \(\left\{f_{n}(x)\right\}\) 的极限函数，并证明 \(\left\{f_{n}(x)\right\}\) 在 \((-\infty,+\infty)\) 上一致收敛．

\subsection{题目7}
七．讨论 \(\int_{1}^{+\infty} \frac{\sin x \cos \frac{1}{x}}{x} \mathrm{~d} x\) 是绝对收敛，条件收敛还是发散？

\subsection{题目8}
八．点 \(P\left(x_{0}, y_{0}, z_{0}\right)\) 是球面 \(x^{2}+y^{2}+z^{2}=1\) 上的点 \(\left(x_{0}, y_{0}, z_{0}>0\right)\) ，求球面在点 \(P\) 的切平面与三个坐标平面围成的四面体体积的最小值。

\subsection{题目9}
九．求 \(f(x)=e^{\sqrt{x}}\) 在 \([0,2 \pi]\) 上的傅里叶级数，并求级数的收敛函数以及 \(\sum_{n=1}^{\infty} \frac{1}{n^{2}+2}\) 的值．

附加题．\(f(x, y)\) 在 \((0,1)\) 附近存在一阶连续的偏导数，且 \(f_{y}(0,1) \neq 0, f(0,1)=0\) 。证明：\(f\left(x, \int_{0}^{t} \sin u \mathrm{~d} u\right)=0\)\\
在点 \(\left(0, \frac{\pi}{2}\right)\) 附近唯一确定单值函数 \(t=\varphi(x)\) ，并求 \(\varphi^{\prime}(0)\) ．

\subsection{题目10}
十一．椭球面 \(S: \frac{x^{2}}{a^{2}}+\frac{y^{2}}{b^{2}}+\frac{z^{2}}{c^{2}}=1(z \geq 0),(\cos \alpha, \cos \beta, \cos \gamma)\) 是椭球面上的单位外法向量，求曲面积分

\[
\iint_{S} z\left(\frac{x}{a^{2}} \cos \alpha+\frac{y}{b^{2}} \cos \beta+\frac{z}{c^{2}} \cos \gamma\right) \mathrm{d} S
\]

\subsection{题目11}
十二．求重积分 \(\iint_{D} \frac{(\sqrt{x}+\sqrt{y})^{4}}{x^{2}} \mathrm{~d} x \mathrm{~d} y\) ，其中 \(D\) 为 \(\sqrt{x}+\sqrt{y}=1, \sqrt{x}+\sqrt{y}=2\) 与 \(y=x\) 及 \(x\) 轴所围区域．
\end{document}
